\textbf{Functional requirements:}
\begin{itemize}
	\item Reconstruct a 3D point cloud of the sea surface in real time (per recorded frame).
	\item Georeference that point cloud by fusing it with an independent navigation/pose solution.
	\item Estimate sea-state parameters (significant wave height ($H_s$), peak period ($T_p$) and, where feasible, propagation direction).
	\item Validate the point cloud by cross-checking the sea-state parameters against an independent sensing modality.
\end{itemize}

\textbf{Non-functional requirements:}
\begin{itemize}
	\item The system must operate with minimum latency, with a target of $\leq$\, 200ms from frame capture to point-cloud output.
	\item The system must be reliable and robust to the specific challenges of the sea surface (low texture, specularity) and to the operational conditions of a moving platform (vibrations, rolling/pitching, changing lighting conditions).
	\item The system must be modular and maintainable, allowing for easy integration of new sensors or algorithms, and for future upgrades to the stereo-matching backend or wave-estimation methods.
\end{itemize}

\textbf{Constraints:}
\begin{itemize}
	\item \begin{sloppypar}The system must operate on a single embedded edge-AI module (NVIDIA Jetson Orin Nano-class, 8~GB shared memory), which imposes strict constraints on the GPU budget available for disparity estimation.\end{sloppypar}
	\item Development and testing were limited to the hardware and tooling available: a single development workstation, the existing test structure, with limited scope to alter the physical setup, in order to mimic the conditions of the final deployment.
	\item Validation depended on recorded datasets, as due to time constraints and meteorological conditions, no further test campaigns could be run in order to gather data under different or rougher sea states and test the pipeline live on the vessel.
\end{itemize}
